\section{Codebook for the Open Questionnaire Items}
\label{annex:qualitative-codebook}

This annex documents the codebook used to analyze the free-text answers to
items F1 and F2 (\autoref{annex:davison-questionnaire}).
F1 asks for a situation in which the agent intervened and for its effect on the
discussion.
F2 asks what was helpful, what was disruptive, and what was missing.
\autoref{tab:qualitative-codebook} lists the codes.

Each coding has three levels: a main category, a subcategory, and a polarity
(negative, positive, or neutral).
Polarity is an attribute of the coding, not a code of its own.
Anchor examples keep the original German wording, including the formulations as
written.
Names stay anonymized as \texttt{\{Name von Gx-Px\}}.
One passage may anchor more than one code.
D1 and D2 belong to context understanding.
D3 is filed with the E~codes under acceptance and group dynamics, because it
describes the group's reaction rather than a failure of context tracking.
W1 is a meta-code: every explicit suggestion receives W1 together with the
topical code named in that row.

\subsection*{General coding rules}

\begin{itemize}
  \item \textbf{Multiple coding.}
    One segment can receive several codes.
    The marker \enquote{\ldots} separates related sub-statements within an answer;
    each sub-statement is coded on its own.
    Every explicit improvement suggestion additionally receives W1.
    Valence is recorded as polarity, not as a separate code.
  \item \textbf{Explicitness.}
    Code only what is stated explicitly or by a clear implication.
    Do not interpret unclear statements.
    Passages flagged \texttt{[!sic]} in the material are not coded.
  \item \textbf{Duplicates and anonymization.}
    Duplicate quotes are coded as separate units.
    Names in curly brackets remain anonymized.
  \item \textbf{Attributes recorded with each coding.}
    Group (G1--G3), person (P1--P4), and intensity (strong, medium, or weak).
    Intensity is assigned only when the wording states it explicitly.
\end{itemize}

\begingroup
  \footnotesize
  \setlength{\tabcolsep}{4pt}
  \setlength{\extrarowheight}{0.4pt}
  \renewcommand{\arraystretch}{1.05}
  \setlength{\LTpre}{0.4em}
  \setlength{\LTpost}{0.6em}
  \setlength{\emergencystretch}{2em}
  \newcommand{\anon}[1]{\{#1\}}
  \newcommand{\cbgap}{\addlinespace[0.5em]}
  \newcommand{\cbcatbreak}{%
    \addlinespace[0.15em]\midrule\addlinespace[0.4em]}
  \newcommand{\cbcat}[1]{%
    \multicolumn{3}{@{}l@{}}{\textbf{#1}} \\*[0.2em]}
  \newcommand{\cbrow}[6]{%
    \textbf{#1}\par\vspace{0.12em}#2\par\vspace{0.1em}{\itshape #3}
    & #4\par\vspace{0.32em}\textit{Coding rule:} #5
    & \enquote{#6} \\
  }
  \begin{xltabular}{\textwidth}{@{}
      >{\RaggedRight\arraybackslash}p{3.55cm}
      >{\RaggedRight\arraybackslash}X
      >{\RaggedRight\arraybackslash}X
    @{}}
    \caption{Codebook for the free-text answers to questionnaire items F1 and F2.
      Anchor examples are quoted in the original German.}
    \label{tab:qualitative-codebook}\\[0.7em]
    \toprule
    \textbf{Code and subcategory} & \textbf{Definition} & \textbf{Anchor example} \\
    \midrule
    \endfirsthead
    \caption[]{Codebook for the free-text answers to questionnaire items F1 and F2
      (continued).}\\[0.7em]
    \toprule
    \textbf{Code and subcategory} & \textbf{Definition} & \textbf{Anchor example} \\
    \midrule
    \endhead
    \bottomrule
    \endfoot
    \bottomrule
    \endlastfoot
    \cbcat{Intervention Behavior}
    \cbrow
      {A1}
      {Interruption of conversation flow}
      {negative}
      {The AI interjects during a human contribution, interrupts a speaker, or its
      behavior is perceived as disrupting the conversation flow.}
      {Code if interruption or flow disruption is explicitly addressed; disruption
      without interruption is included; multiple coding with A3 is possible.}
      {Sie ist nach jedem Redebeitrag angesprungen… … und hat den Redefluss dann
      sehr stark eingeschränkt.}
    \cbgap
    \cbrow
      {A1}
      {Interruption of conversation flow}
      {positive}
      {The AI interrupts in order to create space for others (helpful interruption).}
      {Code when the interruption serves a valued purpose, even if it is still
      experienced as disturbing.}
      {Positiv: In meinen Augen war hilfreich das sie mich unterbrochen hat um
      \anon{Name von G3-P1} und \anon{Name von G3-P4} mehr raum zu geben, das war
      gut, auch wenn es mich gestört hat :D}
    \cbgap
    \cbrow
      {A2}
      {Intervention frequency too high}
      {negative}
      {Statement that the AI intervenes or interrupts too often.}
      {Focus on frequency, not timing; \enquote{should only intervene selectively}
      is also coded here.}
      {Negativ: Die KI hat viel zu oft unterbrocheniner Meinung nach sollte die KI
      nur punktuell eingreifen wenn das Gespräch wirklich abdriftet oder wenn nur
      wenige Personen den ganzen Redeanteil haben.}
    \cbgap
    \cbrow
      {A2}
      {Intervention frequency too low}
      {negative}
      {Statement that the AI intervenes too rarely (for example, when nobody knows
      how to proceed).}
      {Only code if a lack of intervention at needed moments is explicit.}
      {Auf der anderen Seite habe ich auch zwei Situationen im Kopf, indem die
      Diskussion gestoppt hat wer nicht mehr genau wussten, wer jetzt etwas sagen
      soll, wie es weitergeht und an diesen Stellen hat sie nicht eingegriffen.}
    \cbgap
    \cbrow
      {A3}
      {Timing / delay}
      {negative}
      {The AI reacts with a delay, too late, or lags behind the group's thinking
      process.}
      {Focus on delay, not interruption; combination with A1 is possible; slow
      speech tempo is also coded here, in addition to C2.}
      {Inhaltlich waren die Beiträge gut,… … allerdings verzögert, oder hinter
      unserem Denkprozess hinterher,…}
    \cbgap
    \cbrow
      {A4}
      {Inappropriate moment / disregard of context}
      {negative}
      {The AI intervenes at substantively inappropriate moments (for example, after
      consensus was reached) or re-addresses settled points.}
      {Distinguish from A3 (purely temporal): here the substantive fit of the moment
      matters; re-addressing settled content also receives D2.}
      {Man ist mitten in einem Diskurs, und wird von einem Einwand der KI
      unterbrochen, … … wo man bereits zu einem Konsens sich geeinigt hat.}
    \cbcatbreak
    \cbcat{Moderation Behavior}
    \cbrow
      {B1}
      {Summarizing the last contribution}
      {negative}
      {The AI summarizes the previous contribution before its own statement;
      experienced as disrupting the meeting flow.}
      {For short meetings mostly negative; split the polarity if both evaluations
      appear.}
      {Negativ: Die KI hat nach dem Bedanken… … oft das zuletzt gesagte
      zusammengefasst, das hat den Meetingfluss negativ beeinflusst.}
    \cbgap
    \cbrow
      {B1}
      {Summarizing the last contribution}
      {positive}
      {Summarizing is experienced as helpful, for example when earlier content is
      recalled in a long meeting.}
      {Code when usefulness is explicitly stated.}
      {… Das könnte hilfreich sein, wenn man längere Meetings hat und nochmal
      Inhalte vom Anfang zusammengefasst bekommt.}
    \cbgap
    \cbrow
      {B2}
      {Standardized phrases (thanks, formulas)}
      {negative}
      {The AI repeatedly uses identical rhetorical patterns
      (\enquote{thank you xy}), experienced as unnecessary or time-consuming.}
      {Code only if repetition or standardization is explicitly addressed.}
      {Negativ: Die KI hat jedes Mal \enquote{Danke xy} gesagt, was nach dem zweiten
      Mal nicht mehr nötig war.}
    \cbgap
    \cbrow
      {B3}
      {Including quiet participants / balancing speaking time}
      {positive}
      {The AI explicitly addresses participants who have said little, in order to
      balance speaking time and hear all perspectives.}
      {Often positive; cases in which the goal fails (a person is asked although
      already speaking) are also coded here, but as negative; combine with A1 if
      the turn is felt as an interruption.}
      {Positiv: \anon{Name von G1-P4} ist eher stiller bei der Diskussion gewesen
      und wurde durch die KI gezielt angesprochen, sodass auch seine Meinung in
      Ruhe angehört werden konnte.}
    \cbgap
    \cbrow
      {B3}
      {Including quiet participants / balancing speaking time}
      {negative}
      {The inclusion attempt misfires (for example, a person is asked although
      already saying exactly that, or not all members are asked).}
      {Code when the mechanism fails or is experienced as disruptive despite the
      positive goal.}
      {Die KI wollte, dass \anon{Name von G3-P4} ihren Beitrag sagt, obwohl sie
      schon angefangen hatte genau das zu sagen, was die KI gefragt hat.}
    \cbgap
    \cbrow
      {B4}
      {Content-related impulses / steering}
      {positive}
      {The AI provides helpful impulses, steers the direction of the conversation,
      or gives recognition.}
      {Distinguish from steering toward an own agenda; keep polarity separate per
      segment.}
      {Positiv: Hilfreich war das ab und an Impulse gesetzt wurden, um den
      Gesprächsfluss weiter in die richtige Richtung zu lenken und… … auch um
      Anerkennung zu zollen}
    \cbgap
    \cbrow
      {B4}
      {Structuring / common thread}
      {negative}
      {A lack of a guiding structure or framing for the meeting is criticized.}
      {Only code when the missing common thread is explicit.}
      {Negativ: Gefehlt hat ein besserer inhaltlicher Bezug mit sinnvollem Timing,..
      … sowie eine rote Linie, die dem Meeting einen Ramen gibt.}
    \cbgap
    \cbrow
      {B4}
      {Own agenda / external steering}
      {negative}
      {The AI pursues its own agenda or steers against the approach the group has
      agreed.}
      {Code when the AI acts against the group consensus or goal.}
      {Negativ: Störend war, dass die KI gefühlt versucht hat eine eigene Agenda zu
      verfolgen bzw. … … sich nicht genug dem in der Gruppe abgestimmten Vorgehen
      angepasst hat.}
    \cbgap
    \cbrow
      {B5}
      {Resolving misunderstandings / added value}
      {positive}
      {The AI recognizes and resolves content-related misunderstandings or divergent
      states of information.}
      {Only explicit examples of a substantively helpful intervention.}
      {Positiv: Die Situation, wo die KI aufgelöst hat, dass wir unterschiedliche
      Informationen hatten, war wahrscheinlich das meist hilfreichste.}
    \cbcatbreak
    \cbcat{Delivery}
    \cbrow
      {C1}
      {Contribution length / precision}
      {negative}
      {Statements about the length or precision of AI contributions (too long,
      talked too much, a shorter contribution was desired).}
      {Combine with A3 when speed is also criticized.}
      {Die Antworten könnten noch etwas kürzer, präziser und schneller sein, weil
      die KI hat schon sehr viel geredet.}
    \cbgap
    \cbrow
      {C2}
      {Speech/output characteristics (tempo, animation, sound)}
      {negative}
      {Articulation and speech tempo are perceived as too slow or disruptive.}
      {Slow speech tempo that disrupts the flow also receives A1.}
      {Negativ: Das Sprachtempo der KI war zu langsam. Hat den Gesprächsfluss
      gestört.}
    \cbgap
    \cbrow
      {C2}
      {Speech/output characteristics (tempo, animation, sound)}
      {positive}
      {Animation and sound at the onset of speech are experienced as appealing.}
      {Code when the audiovisual presentation is explicitly praised.}
      {Die Animation und der sound waren cool, wenn sie [die KI] angefangen hat zu
      sprechen.}
    \cbgap
    \cbrow
      {C3}
      {Aborted/restarted utterances}
      {negative}
      {The AI aborts utterances mid-sentence and starts anew; multiple restarts.}
      {Only the AI's own behavior, not a human interrupting the AI.}
      {Außerdem hat die KI oft angefangen zu reden und mehrfach im Satz noch einmal
      erneut angefangen zu reden mit einem neuen Satz.}
    \cbgap
    \cbrow
      {C4}
      {Parallel speech / intelligibility}
      {negative}
      {The AI keeps talking while humans continue speaking, so that the actual
      speaker is hard to understand.}
      {Combine with A1, since the flow is also disturbed.}
      {Und durch das penetrante weiterreden der KI,… … während andere Leute dann
      auch nicht aufgehört haben zu reden, war es super schwer die eigentliche
      Person zu verstehen.}
    \cbcatbreak
    \cbcat{Context Understanding}
    \cbrow
      {D1}
      {Confusion of persons/candidates}
      {negative}
      {The AI confuses candidates or participants.}
      {Factual errors are also coded here; combination with D2 is possible.}
      {Oft hat sie die Kandidaten verwechselt und… … war nicht präzise auf dem
      aktuellen Stand der Diskussion, worüber wir gerade geredet haben.}
    \cbgap
    \cbrow
      {D1}
      {Factually false statement}
      {negative}
      {The AI makes factually wrong statements about the content.}
      {Code only concrete false statements about the case content.}
      {Negativ: Die KI hat einmal eine Falschaussage getroffen (Kandidat 3 trinkt
      war ihre aussage)}
    \cbgap
    \cbrow
      {D2}
      {Discussion state not tracked / redundancy}
      {negative}
      {The AI is not up to date; it asks again about things already said or
      processes contributions incorrectly.}
      {Distinguish from D1: here, failing to process content that was already stated
      is the focus.}
      {Negativ: Bestimmte Aspekte wurden mehrfach gefragt oder nicht richtig von der
      KI verstanden/verarbeitet.}
    \cbgap
    \cbrow
      {D2}
      {Turn-taking not recognized}
      {negative}
      {The AI does not register whose turn it is or whether someone has finished.}
      {Code when a failure to track speakers is explicit.}
      {Negativ: Hat nicht gecheckt, wer dran ist und wann er fertig war…}
    \cbgap
    \cbrow
      {D2}
      {Traceability missing (no written contribution)}
      {negative}
      {No written chat contribution exists from which the AI's objection can be
      re-read.}
      {Code the missing traceability of AI statements.}
      {Gefehlt hat außerdem ein schriftlicher Beitrag im Chat, sodass man den
      Einwand nachlesen kann.}
    \cbcatbreak
    \cbcat{Acceptance \& Group Dynamics}
    \cbrow
      {D3}
      {Group reaction: ignoring the AI}
      {negative}
      {The group ignores the AI in reaction to its errors or disruptions.}
      {Code the group's coping strategy, not the triggering error itself.}
      {Dadurch kam es dazu, dass wir als menschliche Meeting Teilnehmer eher dazu
      genannt haben, die ki zu ignorieren, und nicht darauf zu hören, was sie
      sagt.}
    \cbgap
    \cbrow
      {E1}
      {AI as an additional \enquote{party} / disruptive factor}
      {negative}
      {The AI is perceived as an independent party meddling in the social setting of
      the meeting.}
      {Social-perception statements about the AI's role, not single incidents.}
      {Ich denke der Einfluss in einem Teammeeting mit unterschiedlichen Parteien
      und da dann noch die KI ist, was das ganze störend wirken lässt, weil es wie
      noch eine Partei wirkt, die sich einmischt}
    \cbgap
    \cbrow
      {E2}
      {Security and motivation through presence}
      {positive}
      {The mere presence of the AI creates security or motivation to progress
      faster.}
      {Presence effect, not the AI's contributions themselves.}
      {Positiv: … die Anwesenheit von ihr jedoch etwas Sicherheit und Motivation
      gegeben hat, schnell voran zu schreiten.}
    \cbgap
    \cbrow
      {E2}
      {Acceleration through observation pressure}
      {neutral}
      {The meeting accelerates because participants want to avoid the AI stepping in
      again.}
      {Code when the effect is attributed to awareness of the AI, not to its
      suggestions.}
      {Das Meeting hat sich schon sehr schnell und effizient angefüllt, wobei ich
      eher sagen würde das lag nicht an den KI Vorschlägen, sondern an dem
      Bewusstsein, dass wir immer weiter und Progress machen müssen, weil sonst
      sich die ki wieder einmischt und zum weitermachen anregt.}
    \cbgap
    \cbrow
      {E3}
      {Habituation over time}
      {neutral}
      {Initially confusing or disruptive; over time, habituation and better
      collaboration.}
      {Code a temporal change in the experience of the AI.}
      {am Anfang war es erst sehr verwirrend und hat gestört. Man musste erst mal
      reinkommen und mit der KI als Meeting Partner klarkommen. Mit der Zeit wurde
      es dann besser und hat sich dran gewöhnt.}
    \cbgap
    \cbrow
      {E4}
      {Prototype status / maturity}
      {negative}
      {General assessment of technical maturity: a commendable goal, but not yet
      useful.}
      {Overall verdicts on the system's readiness, not single incidents.}
      {Ich sehe was die KI erreichen wollte, lobenswertes Ziel alle zur Sprache
      kommen zu lassen und alle Informationen einzubeziehen, … … aber noch zu sehr
      im Prototypenstatus als dass es aktuell nützlich ist.}
    \cbcatbreak
    \cbcat{Improvement Suggestion}
    \cbrow
      {W1}
      {Control / stop button}
      {neutral}
      {Feature idea: a control to abort or stop the AI.}
      {Meta-code: every explicit suggestion receives W1 plus the matching topical
      subcode; here additionally A1.}
      {Eventuell würde ein \enquote{stopp} button, der einfach die KI unterbricht
      helfen das zu verhindern}
    \cbgap
    \cbrow
      {W1}
      {More subtle input / visual hints}
      {neutral}
      {Feature idea: proactive hints shown only visually to the current speaker.}
      {Additionally A1 and B4.}
      {Ein subtilerer Ansatz mit hinweisen, die proaktiv aber vielleicht nur visuell
      als ein Art \enquote{Hints} nur dem Sprecher angezeigt werden, wäre vielleicht
      mal interessant zu probieren.}
    \cbgap
    \cbrow
      {W1}
      {Written chat contribution}
      {neutral}
      {Feature idea: written AI contributions for traceability.}
      {Additionally D2.}
      {Gefehlt hat außerdem ein schriftlicher Beitrag im Chat, sodass man den
      Einwand nachlesen kann.}
    \cbgap
    \cbrow
      {W1}
      {Introduction / framing at the start}
      {neutral}
      {Feature idea: an AI introduction that sets the framing at the beginning of
      the meeting.}
      {Additionally B4.}
      {Und zuletzt hätte eine Introduction, der ki noch gefehlt vielleicht der ein
      bisschen das Framing der Diskussion am Anfang setzt und man wird nicht
      einfach blind reingeschmissen wird und muss plötzlich anfangen zu
      diskutieren und zu reden.}
    \cbgap
    \cbrow
      {W1}
      {Time / progress cues}
      {neutral}
      {Feature idea: time or progress hints, especially toward the end of the
      meeting.}
      {Additionally A2 and B4.}
      {Es hat für mich ein Hinweis auf die Ablaufende Zeit als Moderator gefehlt,
      hier hatte ich erwartet dass sie gegen Ende mehr eingreift.}
    \cbgap
    \cbrow
      {W1}
      {Individual assistant instead of group AI}
      {neutral}
      {Feature idea: a personal assistant for each participant instead of one group
      AI.}
      {Additionally E1.}
      {… ich denke, wenn da jeder seinen ganz persönlichen Assistenten hätte, wäre
      das viel schneller akzeptiert, weil es weniger aufbringt sondern als support
      auch wahrgenommen werden kann.}
    \cbgap
    \cbrow
      {W1}
      {Selective intervention}
      {neutral}
      {Recommendation: the AI should intervene only selectively (the conversation
      drifts, or speaking time is unbalanced).}
      {Additionally A2.}
      {Meiner Meinung nach sollte die KI nur punktuell eingreifen wenn das Gespräch
      wirklich abdriftet oder wenn nur wenige Personen den ganzen Redeanteil
      haben.}
    \cbgap
    \cbrow
      {W1}
      {Content accuracy / up-to-dateness}
      {neutral}
      {Recommendation: correct references to candidates and up-to-date tracking of
      the context.}
      {Additionally D1 and D2.}
      {Negativ: Gefehlt hat ein besserer inhaltlicher Bezug mit sinnvollem Timing,..}
  \end{xltabular}
\endgroup
